% =============================================================
%  SPURIOUS CORRELATION
%  % =============================================================
%  SPURIOUS CORRELATION
%  % =============================================================
%  SPURIOUS CORRELATION
%  % =============================================================
%  SPURIOUS CORRELATION
%  \input{slides/Spurious_correlation}
% =============================================================


% ------------------------------------------------------------------
%  SLIDE 1 – Spurious Correlation (Pearson)
% ------------------------------------------------------------------

\begin{frame}{Correlation}
  \note{
Now, I extend this example to 500 different job opportunities.

In other words, I generated 500 similar datasets, although only a few rows are shown in this table. Each row represents the number of applicants for one job opportunity.

If we calculate the correlation between applicants with a Mathematics background and applicants with a Statistics background in the original count data, the result is close to zero, as you can see.

This is exactly what we expect, because the variables were generated independently.

Now, if we transform the same data into the simplex space, we obtain a very different result. The Mathematics and Statistics components become strongly negatively correlated.

So, what does this correlation actually mean? What kind of information can we extract from it?

This is not an easy question to answer.

A common mistake is to interpret this as an ordinary correlation between applicants from Mathematics and Statistics backgrounds, as we would do in a standard data table.

However, this interpretation is misleading.

These types of correlations are known as spurious correlations.

      }
\only<1>{
  \begin{columns}
    \column{0.5\textwidth}
    \textbf{Original Data (Counts)}

    \vspace{0.2cm}
    \begin{center}
    \begin{tabular}{ccc}
      \hline
      Math & Stats & CS \\
      \hline
      48  & 230 & 59  \\
      286 & 181 & 73  \\
      286 & 98  & 147 \\
      35  & 14  & 12  \\
      255 & 104 & 235 \\
      $\vdots$ & $\vdots$ & $\vdots$ \\
      \hline
    \end{tabular}
    \end{center}

    \vspace{0.2cm}
    \[
      \mathrm{corr}(\text{Math},\,\text{Stats}) = 0.087
    \]

    \column{0.5\textwidth}
    
  \end{columns}


}
\only<2>{
\begin{columns}
    \column{0.5\textwidth}
    \textbf{Original Data (Counts)}

    \vspace{0.2cm}
    \begin{center}
    \begin{tabular}{ccc}
      \hline
      Math & Stats & CS \\
      \hline
      48  & 230 & 59  \\
      286 & 181 & 73  \\
      286 & 98  & 147 \\
      35  & 14  & 12  \\
      255 & 104 & 235 \\
      $\vdots$ & $\vdots$ & $\vdots$ \\
      \hline
    \end{tabular}
    \end{center}

    \vspace{0.2cm}
    \[
      \mathrm{corr}(\text{Math},\,\text{Stats}) = 0.087
    \]

    \column{0.5\textwidth}
    \textbf{Compositional Data (Normalized to 100)}

    \vspace{0.2cm}
    \begin{center}
    \begin{tabular}{ccc}
      \hline
      Math & Stats & CS \\
      \hline
      14.2 & 68.2 & 17.5 \\
      53.0 & 33.5 & 13.5 \\
      53.9 & 18.5 & 27.7 \\
      57.4 & 23.0 & 19.7 \\
      42.9 & 17.5 & 39.6 \\
      $\vdots$ & $\vdots$ & $\vdots$ \\
      \hline
    \end{tabular}
    \end{center}

    \vspace{0.2cm}
    \[
      \mathrm{corr}(X_{\text{math}},\, X_{\text{stats}}) = -0.874
    \]
  \end{columns}

}

\end{frame}

\end{frame}
\begin{frame}{Spurious Correlation}

 


\only<1>{
\note{Remove this slide:
Spurious correlation

One important consequence of working with compositional data is spurious correlation.

This phenomenon was introduced by Pearson in 1897.

He showed that ratios that share a common denominator can appear correlated, even when the original variables are completely independent.

To see this, let us consider three independent random variables: A, B, and C.

Now suppose we construct two new variables by dividing A by C and B by C.

Even though A and B are independent, these two ratios can become correlated because they share the same denominator.

}
  \begin{block}{Pearson (1897)}
    Ratios that share a common denominator are correlated even when the
    original variables are \emph{completely independent}.
  \end{block}

  \vspace{0.4em}
  \textbf{Simple illustration:} Let $A$, $B$, $C$ be mutually independent. Then:
  \[
    \operatorname{corr}\!\left(\frac{A}{C},\;\frac{B}{C}\right) \neq 0
  \]
  \vspace{0.2em}
  The shared denominator $C$ couples the two ratios artificially.
  The observed correlation reflects the \textbf{denominator's variability},
  not any true relationship between $A$ and $B$.
}

\only<2>{
\note{Spurious correlation

One important consequence of working with compositional data is spurious correlation.

This phenomenon was introduced by Pearson in 1897.

He showed that ratios that share a common denominator can appear correlated, even when the original variables are completely independent.

To see this, let us consider three independent random variables: A, B, and C.

Now suppose we construct two new variables by dividing A by C and B by C.

Even though A and B are independent, these two ratios can become correlated because they share the same denominator.

We observe exactly the same phenomenon in compositional data.

Since all components are normalized by the same sum, a common dependency is introduced.

As a result, even if two components are independent in the original space, they may appear dependent in the simplex space.

However, this is not a mistake.

}
  \begin{block}{Pearson (1897)}
    Ratios that share a common denominator are correlated even when the
    original variables are \emph{completely independent}.
  \end{block}

  \vspace{0.4em}

  \begin{columns}[t]
    \begin{column}{0.38\textwidth}
      \textbf{Original vector:}\\[4pt]
      $\mathbf{Y} = (Y_1, Y_2, \ldots, Y_K) \in \mathbb{R}_+^K$
    \end{column}
    \begin{column}{0.58\textwidth}
      \textbf{Simplex (normalized):}\\[4pt]
      $\mathbf{s}(\mathbf{Y}) = \!\left(\dfrac{Y_1}{S},\, \dfrac{Y_2}{S},\, \ldots,\, \dfrac{Y_K}{S}\right),
      \quad S = \sum_{j=1}^{K} Y_j$
    \end{column}
  \end{columns}

  \vspace{0.5em}
  \rule{\textwidth}{0.3pt}
  \vspace{0.3em}

  Even if $\operatorname{corr}(Y_1,\, Y_2) = 0$, we usually observe:
  \[
    \boxed{\operatorname{corr}\!\left(\frac{Y_1}{S},\,\frac{Y_2}{S}\right) \neq 0}
  \]
}

\only<3>{
\note{Mathematically, everything is correct.

The problem is only interpretation.

Pearson’s warning was exactly about this point: we should not interpret correlations between ratios as correlations between the original variables.

Unfortunately, this mistake often happens in practice.

}
  \[
  \boxed{
    \operatorname{corr}\!\left(\frac{Y_1}{S}, \frac{Y_2}{S}\right)
    =
    \frac{
      \mathbb{E}\!\left(\frac{Y_1 Y_2}{S^2}\right)
      - \mathbb{E}\!\left(\frac{Y_1}{S}\right)\mathbb{E}\!\left(\frac{Y_2}{S}\right)
    }{
      \sqrt{\operatorname{Var}\!\left(\frac{Y_1}{S}\right)}\;
      \sqrt{\operatorname{Var}\!\left(\frac{Y_2}{S}\right)}
    }
  }
  \]
  \vspace{0.4em}
  \begin{block}{Pearson's warning}
    The danger arises when one interprets a correlation between ratios as evidence
    of a genuine association between the \emph{original} variables.
    Such an inference is \textbf{not justified} -- the correlation may exist
    entirely because of the shared denominator, with no signal from the data at all.
  \end{block}
}



\only<4>{
  \note{
  For example, consider a job application dataset with different academic backgrounds such as mathematics, statistics, and computer science.

In the original count data (On the left), the correlation between mathematics and statistics may be very small, close to zero.

However, after normalization to percentages — that is, after moving to the simplex space — the same variables can become strongly negatively or positively correlated.

      }

  \begin{columns}
    \column{0.5\textwidth}
    \textbf{Original Data (Counts)}

    \vspace{0.2cm}
    \begin{center}
    \begin{tabular}{ccc}
      \hline
      Math & Stats & CS \\
      \hline
      48  & 230 & 59  \\
      286 & 181 & 73  \\
      286 & 98  & 147 \\
      35  & 14  & 12  \\
      255 & 104 & 235 \\
      $\vdots$ & $\vdots$ & $\vdots$ \\
      \hline
    \end{tabular}
    \end{center}

    \vspace{0.2cm}
    \[
      \mathrm{corr}(\text{Math},\,\text{Stats}) = 0.087
    \]

    \column{0.5\textwidth}
    \textbf{Compositional Data (Normalized to 100)}

    \vspace{0.2cm}
    \begin{center}
    \begin{tabular}{ccc}
      \hline
      Math & Stats & CS \\
      \hline
      14.2 & 68.2 & 17.5 \\
      53.0 & 33.5 & 13.5 \\
      53.9 & 18.5 & 27.7 \\
      57.4 & 23.0 & 19.7 \\
      42.9 & 17.5 & 39.6 \\
      $\vdots$ & $\vdots$ & $\vdots$ \\
      \hline
    \end{tabular}
    \end{center}

    \vspace{0.2cm}
    \[
      \mathrm{corr}(X_{\text{math}},\, X_{\text{stats}}) = -0.874
    \]
  \end{columns}

  \vspace{0.3em}
  \textbf{Key message:} Spurious dependence:
  This correlation does \textbf{not} necessarily reflect a true relationship.
}

\end{frame}


% ------------------------------------------------------------------
%  SLIDE 2 – Statistical Difficulties: Singular Matrix
% ------------------------------------------------------------------
\begin{frame}{Statistical Difficulties}

\note{
Singularity of the covariance matrix

Another important consequence is that the covariance matrix is not invertible in the simplex space.

This can be shown quite easily.

Because of the constant-sum constraint, the components are not independent.

As a result, one degree of freedom is lost, and the covariance matrix becomes singular.

This creates a practical limitation.

Many standard statistical methods require the inverse of the covariance matrix, such as regression models and multivariate techniques.

Therefore, these methods cannot be directly applied to compositional data.

There are also several other difficulties that arise from working in the simplex space.

I will not go into all of them here.

However, all of these challenges motivated researchers to develop new methods specifically designed for compositional data.
}

\only<1>{
  \vspace{0.4em}
  \[
  \Sigma_{\mathbf{X}} =
  \begin{pmatrix}
    \operatorname{Var}(X_1)     & \operatorname{cov}(X_1,X_2) & \cdots & \operatorname{cov}(X_1,X_K) \\
    \operatorname{cov}(X_2,X_1) & \operatorname{Var}(X_2)     & \cdots & \operatorname{cov}(X_2,X_K) \\
    \vdots                      & \vdots                       & \ddots & \vdots \\
    \operatorname{cov}(X_K,X_1) & \operatorname{cov}(X_K,X_2) & \cdots & \operatorname{Var}(X_K)
  \end{pmatrix}
  \in \mathbb{R}^{K \times K}
  \]
  \vspace{0.2em}
  The \textbf{last row (and column) is linearly determined} by the others via the
  sum-to-constant constraint, confirming $\operatorname{rank}(\Sigma_{\mathbf{X}}) = K-1$.
}

\only<2>{
  The sum-to-constant constraint causes the \textbf{sample covariance matrix to be singular}:
  \[
    \sum_{i=1}^{K} \operatorname{cov}(X_1, X_i)
    = \operatorname{cov}\!\left(X_1,\, \sum_{i=1}^K X_i\right)
    = \operatorname{cov}(X_1,\, 1) = 0
  \]
  \[
    \Rightarrow\quad
    \sum_{i=2}^{K} \operatorname{cov}(X_1, X_i) = -\operatorname{Var}(X_1).
  \]

  \vspace{0.3em}
  $\Sigma_{\mathbf{X}}$ has rank $K-1$ $\;\Rightarrow\;$ \textbf{not invertible}.

  \vspace{0.5em}
  \begin{columns}[T]
    \begin{column}{0.48\textwidth}
      \textbf{Technical consequence:}
      \begin{itemize}
        \item Regression, LDA, and standard PCA all require $\Sigma^{-1}$ -- unavailable.
      \end{itemize}
    \end{column}
    \begin{column}{0.48\textwidth}
      \textbf{Conceptual consequence:}
      \begin{itemize}
        \item No universally accepted definition of correlation for compositional data exists yet.
      \end{itemize}
    \end{column}
  \end{columns}
}

\end{frame}


% ------------------------------------------------------------------
%  SLIDE 3 – The Log-Ratio Idea
% ------------------------------------------------------------------

% =============================================================


% ------------------------------------------------------------------
%  SLIDE 1 – Spurious Correlation (Pearson)
% ------------------------------------------------------------------

\begin{frame}{Correlation}
  \note{
Now, I extend this example to 500 different job opportunities.

In other words, I generated 500 similar datasets, although only a few rows are shown in this table. Each row represents the number of applicants for one job opportunity.

If we calculate the correlation between applicants with a Mathematics background and applicants with a Statistics background in the original count data, the result is close to zero, as you can see.

This is exactly what we expect, because the variables were generated independently.

Now, if we transform the same data into the simplex space, we obtain a very different result. The Mathematics and Statistics components become strongly negatively correlated.

So, what does this correlation actually mean? What kind of information can we extract from it?

This is not an easy question to answer.

A common mistake is to interpret this as an ordinary correlation between applicants from Mathematics and Statistics backgrounds, as we would do in a standard data table.

However, this interpretation is misleading.

These types of correlations are known as spurious correlations.

      }
\only<1>{
  \begin{columns}
    \column{0.5\textwidth}
    \textbf{Original Data (Counts)}

    \vspace{0.2cm}
    \begin{center}
    \begin{tabular}{ccc}
      \hline
      Math & Stats & CS \\
      \hline
      48  & 230 & 59  \\
      286 & 181 & 73  \\
      286 & 98  & 147 \\
      35  & 14  & 12  \\
      255 & 104 & 235 \\
      $\vdots$ & $\vdots$ & $\vdots$ \\
      \hline
    \end{tabular}
    \end{center}

    \vspace{0.2cm}
    \[
      \mathrm{corr}(\text{Math},\,\text{Stats}) = 0.087
    \]

    \column{0.5\textwidth}
    
  \end{columns}


}
\only<2>{
\begin{columns}
    \column{0.5\textwidth}
    \textbf{Original Data (Counts)}

    \vspace{0.2cm}
    \begin{center}
    \begin{tabular}{ccc}
      \hline
      Math & Stats & CS \\
      \hline
      48  & 230 & 59  \\
      286 & 181 & 73  \\
      286 & 98  & 147 \\
      35  & 14  & 12  \\
      255 & 104 & 235 \\
      $\vdots$ & $\vdots$ & $\vdots$ \\
      \hline
    \end{tabular}
    \end{center}

    \vspace{0.2cm}
    \[
      \mathrm{corr}(\text{Math},\,\text{Stats}) = 0.087
    \]

    \column{0.5\textwidth}
    \textbf{Compositional Data (Normalized to 100)}

    \vspace{0.2cm}
    \begin{center}
    \begin{tabular}{ccc}
      \hline
      Math & Stats & CS \\
      \hline
      14.2 & 68.2 & 17.5 \\
      53.0 & 33.5 & 13.5 \\
      53.9 & 18.5 & 27.7 \\
      57.4 & 23.0 & 19.7 \\
      42.9 & 17.5 & 39.6 \\
      $\vdots$ & $\vdots$ & $\vdots$ \\
      \hline
    \end{tabular}
    \end{center}

    \vspace{0.2cm}
    \[
      \mathrm{corr}(X_{\text{math}},\, X_{\text{stats}}) = -0.874
    \]
  \end{columns}

}

\end{frame}

\end{frame}
\begin{frame}{Spurious Correlation}

 


\only<1>{
\note{Remove this slide:
Spurious correlation

One important consequence of working with compositional data is spurious correlation.

This phenomenon was introduced by Pearson in 1897.

He showed that ratios that share a common denominator can appear correlated, even when the original variables are completely independent.

To see this, let us consider three independent random variables: A, B, and C.

Now suppose we construct two new variables by dividing A by C and B by C.

Even though A and B are independent, these two ratios can become correlated because they share the same denominator.

}
  \begin{block}{Pearson (1897)}
    Ratios that share a common denominator are correlated even when the
    original variables are \emph{completely independent}.
  \end{block}

  \vspace{0.4em}
  \textbf{Simple illustration:} Let $A$, $B$, $C$ be mutually independent. Then:
  \[
    \operatorname{corr}\!\left(\frac{A}{C},\;\frac{B}{C}\right) \neq 0
  \]
  \vspace{0.2em}
  The shared denominator $C$ couples the two ratios artificially.
  The observed correlation reflects the \textbf{denominator's variability},
  not any true relationship between $A$ and $B$.
}

\only<2>{
\note{Spurious correlation

One important consequence of working with compositional data is spurious correlation.

This phenomenon was introduced by Pearson in 1897.

He showed that ratios that share a common denominator can appear correlated, even when the original variables are completely independent.

To see this, let us consider three independent random variables: A, B, and C.

Now suppose we construct two new variables by dividing A by C and B by C.

Even though A and B are independent, these two ratios can become correlated because they share the same denominator.

We observe exactly the same phenomenon in compositional data.

Since all components are normalized by the same sum, a common dependency is introduced.

As a result, even if two components are independent in the original space, they may appear dependent in the simplex space.

However, this is not a mistake.

}
  \begin{block}{Pearson (1897)}
    Ratios that share a common denominator are correlated even when the
    original variables are \emph{completely independent}.
  \end{block}

  \vspace{0.4em}

  \begin{columns}[t]
    \begin{column}{0.38\textwidth}
      \textbf{Original vector:}\\[4pt]
      $\mathbf{Y} = (Y_1, Y_2, \ldots, Y_K) \in \mathbb{R}_+^K$
    \end{column}
    \begin{column}{0.58\textwidth}
      \textbf{Simplex (normalized):}\\[4pt]
      $\mathbf{s}(\mathbf{Y}) = \!\left(\dfrac{Y_1}{S},\, \dfrac{Y_2}{S},\, \ldots,\, \dfrac{Y_K}{S}\right),
      \quad S = \sum_{j=1}^{K} Y_j$
    \end{column}
  \end{columns}

  \vspace{0.5em}
  \rule{\textwidth}{0.3pt}
  \vspace{0.3em}

  Even if $\operatorname{corr}(Y_1,\, Y_2) = 0$, we usually observe:
  \[
    \boxed{\operatorname{corr}\!\left(\frac{Y_1}{S},\,\frac{Y_2}{S}\right) \neq 0}
  \]
}

\only<3>{
\note{Mathematically, everything is correct.

The problem is only interpretation.

Pearson’s warning was exactly about this point: we should not interpret correlations between ratios as correlations between the original variables.

Unfortunately, this mistake often happens in practice.

}
  \[
  \boxed{
    \operatorname{corr}\!\left(\frac{Y_1}{S}, \frac{Y_2}{S}\right)
    =
    \frac{
      \mathbb{E}\!\left(\frac{Y_1 Y_2}{S^2}\right)
      - \mathbb{E}\!\left(\frac{Y_1}{S}\right)\mathbb{E}\!\left(\frac{Y_2}{S}\right)
    }{
      \sqrt{\operatorname{Var}\!\left(\frac{Y_1}{S}\right)}\;
      \sqrt{\operatorname{Var}\!\left(\frac{Y_2}{S}\right)}
    }
  }
  \]
  \vspace{0.4em}
  \begin{block}{Pearson's warning}
    The danger arises when one interprets a correlation between ratios as evidence
    of a genuine association between the \emph{original} variables.
    Such an inference is \textbf{not justified} -- the correlation may exist
    entirely because of the shared denominator, with no signal from the data at all.
  \end{block}
}



\only<4>{
  \note{
  For example, consider a job application dataset with different academic backgrounds such as mathematics, statistics, and computer science.

In the original count data (On the left), the correlation between mathematics and statistics may be very small, close to zero.

However, after normalization to percentages — that is, after moving to the simplex space — the same variables can become strongly negatively or positively correlated.

      }

  \begin{columns}
    \column{0.5\textwidth}
    \textbf{Original Data (Counts)}

    \vspace{0.2cm}
    \begin{center}
    \begin{tabular}{ccc}
      \hline
      Math & Stats & CS \\
      \hline
      48  & 230 & 59  \\
      286 & 181 & 73  \\
      286 & 98  & 147 \\
      35  & 14  & 12  \\
      255 & 104 & 235 \\
      $\vdots$ & $\vdots$ & $\vdots$ \\
      \hline
    \end{tabular}
    \end{center}

    \vspace{0.2cm}
    \[
      \mathrm{corr}(\text{Math},\,\text{Stats}) = 0.087
    \]

    \column{0.5\textwidth}
    \textbf{Compositional Data (Normalized to 100)}

    \vspace{0.2cm}
    \begin{center}
    \begin{tabular}{ccc}
      \hline
      Math & Stats & CS \\
      \hline
      14.2 & 68.2 & 17.5 \\
      53.0 & 33.5 & 13.5 \\
      53.9 & 18.5 & 27.7 \\
      57.4 & 23.0 & 19.7 \\
      42.9 & 17.5 & 39.6 \\
      $\vdots$ & $\vdots$ & $\vdots$ \\
      \hline
    \end{tabular}
    \end{center}

    \vspace{0.2cm}
    \[
      \mathrm{corr}(X_{\text{math}},\, X_{\text{stats}}) = -0.874
    \]
  \end{columns}

  \vspace{0.3em}
  \textbf{Key message:} Spurious dependence:
  This correlation does \textbf{not} necessarily reflect a true relationship.
}

\end{frame}


% ------------------------------------------------------------------
%  SLIDE 2 – Statistical Difficulties: Singular Matrix
% ------------------------------------------------------------------
\begin{frame}{Statistical Difficulties}

\note{
Singularity of the covariance matrix

Another important consequence is that the covariance matrix is not invertible in the simplex space.

This can be shown quite easily.

Because of the constant-sum constraint, the components are not independent.

As a result, one degree of freedom is lost, and the covariance matrix becomes singular.

This creates a practical limitation.

Many standard statistical methods require the inverse of the covariance matrix, such as regression models and multivariate techniques.

Therefore, these methods cannot be directly applied to compositional data.

There are also several other difficulties that arise from working in the simplex space.

I will not go into all of them here.

However, all of these challenges motivated researchers to develop new methods specifically designed for compositional data.
}

\only<1>{
  \vspace{0.4em}
  \[
  \Sigma_{\mathbf{X}} =
  \begin{pmatrix}
    \operatorname{Var}(X_1)     & \operatorname{cov}(X_1,X_2) & \cdots & \operatorname{cov}(X_1,X_K) \\
    \operatorname{cov}(X_2,X_1) & \operatorname{Var}(X_2)     & \cdots & \operatorname{cov}(X_2,X_K) \\
    \vdots                      & \vdots                       & \ddots & \vdots \\
    \operatorname{cov}(X_K,X_1) & \operatorname{cov}(X_K,X_2) & \cdots & \operatorname{Var}(X_K)
  \end{pmatrix}
  \in \mathbb{R}^{K \times K}
  \]
  \vspace{0.2em}
  The \textbf{last row (and column) is linearly determined} by the others via the
  sum-to-constant constraint, confirming $\operatorname{rank}(\Sigma_{\mathbf{X}}) = K-1$.
}

\only<2>{
  The sum-to-constant constraint causes the \textbf{sample covariance matrix to be singular}:
  \[
    \sum_{i=1}^{K} \operatorname{cov}(X_1, X_i)
    = \operatorname{cov}\!\left(X_1,\, \sum_{i=1}^K X_i\right)
    = \operatorname{cov}(X_1,\, 1) = 0
  \]
  \[
    \Rightarrow\quad
    \sum_{i=2}^{K} \operatorname{cov}(X_1, X_i) = -\operatorname{Var}(X_1).
  \]

  \vspace{0.3em}
  $\Sigma_{\mathbf{X}}$ has rank $K-1$ $\;\Rightarrow\;$ \textbf{not invertible}.

  \vspace{0.5em}
  \begin{columns}[T]
    \begin{column}{0.48\textwidth}
      \textbf{Technical consequence:}
      \begin{itemize}
        \item Regression, LDA, and standard PCA all require $\Sigma^{-1}$ -- unavailable.
      \end{itemize}
    \end{column}
    \begin{column}{0.48\textwidth}
      \textbf{Conceptual consequence:}
      \begin{itemize}
        \item No universally accepted definition of correlation for compositional data exists yet.
      \end{itemize}
    \end{column}
  \end{columns}
}

\end{frame}


% ------------------------------------------------------------------
%  SLIDE 3 – The Log-Ratio Idea
% ------------------------------------------------------------------

% =============================================================


% ------------------------------------------------------------------
%  SLIDE 1 – Spurious Correlation (Pearson)
% ------------------------------------------------------------------

\begin{frame}{Correlation}
  \note{
Now, I extend this example to 500 different job opportunities.

In other words, I generated 500 similar datasets, although only a few rows are shown in this table. Each row represents the number of applicants for one job opportunity.

If we calculate the correlation between applicants with a Mathematics background and applicants with a Statistics background in the original count data, the result is close to zero, as you can see.

This is exactly what we expect, because the variables were generated independently.

Now, if we transform the same data into the simplex space, we obtain a very different result. The Mathematics and Statistics components become strongly negatively correlated.

So, what does this correlation actually mean? What kind of information can we extract from it?

This is not an easy question to answer.

A common mistake is to interpret this as an ordinary correlation between applicants from Mathematics and Statistics backgrounds, as we would do in a standard data table.

However, this interpretation is misleading.

These types of correlations are known as spurious correlations.

      }
\only<1>{
  \begin{columns}
    \column{0.5\textwidth}
    \textbf{Original Data (Counts)}

    \vspace{0.2cm}
    \begin{center}
    \begin{tabular}{ccc}
      \hline
      Math & Stats & CS \\
      \hline
      48  & 230 & 59  \\
      286 & 181 & 73  \\
      286 & 98  & 147 \\
      35  & 14  & 12  \\
      255 & 104 & 235 \\
      $\vdots$ & $\vdots$ & $\vdots$ \\
      \hline
    \end{tabular}
    \end{center}

    \vspace{0.2cm}
    \[
      \mathrm{corr}(\text{Math},\,\text{Stats}) = 0.087
    \]

    \column{0.5\textwidth}
    
  \end{columns}


}
\only<2>{
\begin{columns}
    \column{0.5\textwidth}
    \textbf{Original Data (Counts)}

    \vspace{0.2cm}
    \begin{center}
    \begin{tabular}{ccc}
      \hline
      Math & Stats & CS \\
      \hline
      48  & 230 & 59  \\
      286 & 181 & 73  \\
      286 & 98  & 147 \\
      35  & 14  & 12  \\
      255 & 104 & 235 \\
      $\vdots$ & $\vdots$ & $\vdots$ \\
      \hline
    \end{tabular}
    \end{center}

    \vspace{0.2cm}
    \[
      \mathrm{corr}(\text{Math},\,\text{Stats}) = 0.087
    \]

    \column{0.5\textwidth}
    \textbf{Compositional Data (Normalized to 100)}

    \vspace{0.2cm}
    \begin{center}
    \begin{tabular}{ccc}
      \hline
      Math & Stats & CS \\
      \hline
      14.2 & 68.2 & 17.5 \\
      53.0 & 33.5 & 13.5 \\
      53.9 & 18.5 & 27.7 \\
      57.4 & 23.0 & 19.7 \\
      42.9 & 17.5 & 39.6 \\
      $\vdots$ & $\vdots$ & $\vdots$ \\
      \hline
    \end{tabular}
    \end{center}

    \vspace{0.2cm}
    \[
      \mathrm{corr}(X_{\text{math}},\, X_{\text{stats}}) = -0.874
    \]
  \end{columns}

}

\end{frame}

\end{frame}
\begin{frame}{Spurious Correlation}

 


\only<1>{
\note{Remove this slide:
Spurious correlation

One important consequence of working with compositional data is spurious correlation.

This phenomenon was introduced by Pearson in 1897.

He showed that ratios that share a common denominator can appear correlated, even when the original variables are completely independent.

To see this, let us consider three independent random variables: A, B, and C.

Now suppose we construct two new variables by dividing A by C and B by C.

Even though A and B are independent, these two ratios can become correlated because they share the same denominator.

}
  \begin{block}{Pearson (1897)}
    Ratios that share a common denominator are correlated even when the
    original variables are \emph{completely independent}.
  \end{block}

  \vspace{0.4em}
  \textbf{Simple illustration:} Let $A$, $B$, $C$ be mutually independent. Then:
  \[
    \operatorname{corr}\!\left(\frac{A}{C},\;\frac{B}{C}\right) \neq 0
  \]
  \vspace{0.2em}
  The shared denominator $C$ couples the two ratios artificially.
  The observed correlation reflects the \textbf{denominator's variability},
  not any true relationship between $A$ and $B$.
}

\only<2>{
\note{Spurious correlation

One important consequence of working with compositional data is spurious correlation.

This phenomenon was introduced by Pearson in 1897.

He showed that ratios that share a common denominator can appear correlated, even when the original variables are completely independent.

To see this, let us consider three independent random variables: A, B, and C.

Now suppose we construct two new variables by dividing A by C and B by C.

Even though A and B are independent, these two ratios can become correlated because they share the same denominator.

We observe exactly the same phenomenon in compositional data.

Since all components are normalized by the same sum, a common dependency is introduced.

As a result, even if two components are independent in the original space, they may appear dependent in the simplex space.

However, this is not a mistake.

}
  \begin{block}{Pearson (1897)}
    Ratios that share a common denominator are correlated even when the
    original variables are \emph{completely independent}.
  \end{block}

  \vspace{0.4em}

  \begin{columns}[t]
    \begin{column}{0.38\textwidth}
      \textbf{Original vector:}\\[4pt]
      $\mathbf{Y} = (Y_1, Y_2, \ldots, Y_K) \in \mathbb{R}_+^K$
    \end{column}
    \begin{column}{0.58\textwidth}
      \textbf{Simplex (normalized):}\\[4pt]
      $\mathbf{s}(\mathbf{Y}) = \!\left(\dfrac{Y_1}{S},\, \dfrac{Y_2}{S},\, \ldots,\, \dfrac{Y_K}{S}\right),
      \quad S = \sum_{j=1}^{K} Y_j$
    \end{column}
  \end{columns}

  \vspace{0.5em}
  \rule{\textwidth}{0.3pt}
  \vspace{0.3em}

  Even if $\operatorname{corr}(Y_1,\, Y_2) = 0$, we usually observe:
  \[
    \boxed{\operatorname{corr}\!\left(\frac{Y_1}{S},\,\frac{Y_2}{S}\right) \neq 0}
  \]
}

\only<3>{
\note{Mathematically, everything is correct.

The problem is only interpretation.

Pearson’s warning was exactly about this point: we should not interpret correlations between ratios as correlations between the original variables.

Unfortunately, this mistake often happens in practice.

}
  \[
  \boxed{
    \operatorname{corr}\!\left(\frac{Y_1}{S}, \frac{Y_2}{S}\right)
    =
    \frac{
      \mathbb{E}\!\left(\frac{Y_1 Y_2}{S^2}\right)
      - \mathbb{E}\!\left(\frac{Y_1}{S}\right)\mathbb{E}\!\left(\frac{Y_2}{S}\right)
    }{
      \sqrt{\operatorname{Var}\!\left(\frac{Y_1}{S}\right)}\;
      \sqrt{\operatorname{Var}\!\left(\frac{Y_2}{S}\right)}
    }
  }
  \]
  \vspace{0.4em}
  \begin{block}{Pearson's warning}
    The danger arises when one interprets a correlation between ratios as evidence
    of a genuine association between the \emph{original} variables.
    Such an inference is \textbf{not justified} -- the correlation may exist
    entirely because of the shared denominator, with no signal from the data at all.
  \end{block}
}



\only<4>{
  \note{
  For example, consider a job application dataset with different academic backgrounds such as mathematics, statistics, and computer science.

In the original count data (On the left), the correlation between mathematics and statistics may be very small, close to zero.

However, after normalization to percentages — that is, after moving to the simplex space — the same variables can become strongly negatively or positively correlated.

      }

  \begin{columns}
    \column{0.5\textwidth}
    \textbf{Original Data (Counts)}

    \vspace{0.2cm}
    \begin{center}
    \begin{tabular}{ccc}
      \hline
      Math & Stats & CS \\
      \hline
      48  & 230 & 59  \\
      286 & 181 & 73  \\
      286 & 98  & 147 \\
      35  & 14  & 12  \\
      255 & 104 & 235 \\
      $\vdots$ & $\vdots$ & $\vdots$ \\
      \hline
    \end{tabular}
    \end{center}

    \vspace{0.2cm}
    \[
      \mathrm{corr}(\text{Math},\,\text{Stats}) = 0.087
    \]

    \column{0.5\textwidth}
    \textbf{Compositional Data (Normalized to 100)}

    \vspace{0.2cm}
    \begin{center}
    \begin{tabular}{ccc}
      \hline
      Math & Stats & CS \\
      \hline
      14.2 & 68.2 & 17.5 \\
      53.0 & 33.5 & 13.5 \\
      53.9 & 18.5 & 27.7 \\
      57.4 & 23.0 & 19.7 \\
      42.9 & 17.5 & 39.6 \\
      $\vdots$ & $\vdots$ & $\vdots$ \\
      \hline
    \end{tabular}
    \end{center}

    \vspace{0.2cm}
    \[
      \mathrm{corr}(X_{\text{math}},\, X_{\text{stats}}) = -0.874
    \]
  \end{columns}

  \vspace{0.3em}
  \textbf{Key message:} Spurious dependence:
  This correlation does \textbf{not} necessarily reflect a true relationship.
}

\end{frame}


% ------------------------------------------------------------------
%  SLIDE 2 – Statistical Difficulties: Singular Matrix
% ------------------------------------------------------------------
\begin{frame}{Statistical Difficulties}

\note{
Singularity of the covariance matrix

Another important consequence is that the covariance matrix is not invertible in the simplex space.

This can be shown quite easily.

Because of the constant-sum constraint, the components are not independent.

As a result, one degree of freedom is lost, and the covariance matrix becomes singular.

This creates a practical limitation.

Many standard statistical methods require the inverse of the covariance matrix, such as regression models and multivariate techniques.

Therefore, these methods cannot be directly applied to compositional data.

There are also several other difficulties that arise from working in the simplex space.

I will not go into all of them here.

However, all of these challenges motivated researchers to develop new methods specifically designed for compositional data.
}

\only<1>{
  \vspace{0.4em}
  \[
  \Sigma_{\mathbf{X}} =
  \begin{pmatrix}
    \operatorname{Var}(X_1)     & \operatorname{cov}(X_1,X_2) & \cdots & \operatorname{cov}(X_1,X_K) \\
    \operatorname{cov}(X_2,X_1) & \operatorname{Var}(X_2)     & \cdots & \operatorname{cov}(X_2,X_K) \\
    \vdots                      & \vdots                       & \ddots & \vdots \\
    \operatorname{cov}(X_K,X_1) & \operatorname{cov}(X_K,X_2) & \cdots & \operatorname{Var}(X_K)
  \end{pmatrix}
  \in \mathbb{R}^{K \times K}
  \]
  \vspace{0.2em}
  The \textbf{last row (and column) is linearly determined} by the others via the
  sum-to-constant constraint, confirming $\operatorname{rank}(\Sigma_{\mathbf{X}}) = K-1$.
}

\only<2>{
  The sum-to-constant constraint causes the \textbf{sample covariance matrix to be singular}:
  \[
    \sum_{i=1}^{K} \operatorname{cov}(X_1, X_i)
    = \operatorname{cov}\!\left(X_1,\, \sum_{i=1}^K X_i\right)
    = \operatorname{cov}(X_1,\, 1) = 0
  \]
  \[
    \Rightarrow\quad
    \sum_{i=2}^{K} \operatorname{cov}(X_1, X_i) = -\operatorname{Var}(X_1).
  \]

  \vspace{0.3em}
  $\Sigma_{\mathbf{X}}$ has rank $K-1$ $\;\Rightarrow\;$ \textbf{not invertible}.

  \vspace{0.5em}
  \begin{columns}[T]
    \begin{column}{0.48\textwidth}
      \textbf{Technical consequence:}
      \begin{itemize}
        \item Regression, LDA, and standard PCA all require $\Sigma^{-1}$ -- unavailable.
      \end{itemize}
    \end{column}
    \begin{column}{0.48\textwidth}
      \textbf{Conceptual consequence:}
      \begin{itemize}
        \item No universally accepted definition of correlation for compositional data exists yet.
      \end{itemize}
    \end{column}
  \end{columns}
}

\end{frame}


% ------------------------------------------------------------------
%  SLIDE 3 – The Log-Ratio Idea
% ------------------------------------------------------------------

% =============================================================


% ------------------------------------------------------------------
%  SLIDE 1 – Spurious Correlation (Pearson)
% ------------------------------------------------------------------

\begin{frame}{Correlation}
  \note{
Now, I extend this example to 500 different job opportunities.

In other words, I generated 500 similar datasets, although only a few rows are shown in this table. Each row represents the number of applicants for one job opportunity.

If we calculate the correlation between applicants with a Mathematics background and applicants with a Statistics background in the original count data, the result is close to zero, as you can see.

This is exactly what we expect, because the variables were generated independently.

Now, if we transform the same data into the simplex space, we obtain a very different result. The Mathematics and Statistics components become strongly negatively correlated.

So, what does this correlation actually mean? What kind of information can we extract from it?

This is not an easy question to answer.

A common mistake is to interpret this as an ordinary correlation between applicants from Mathematics and Statistics backgrounds, as we would do in a standard data table.

However, this interpretation is misleading.

These types of correlations are known as spurious correlations.

      }
\only<1>{
  \begin{columns}
    \column{0.5\textwidth}
    \textbf{Original Data (Counts)}

    \vspace{0.2cm}
    \begin{center}
    \begin{tabular}{ccc}
      \hline
      Math & Stats & CS \\
      \hline
      48  & 230 & 59  \\
      286 & 181 & 73  \\
      286 & 98  & 147 \\
      35  & 14  & 12  \\
      255 & 104 & 235 \\
      $\vdots$ & $\vdots$ & $\vdots$ \\
      \hline
    \end{tabular}
    \end{center}

    \vspace{0.2cm}
    \[
      \mathrm{corr}(\text{Math},\,\text{Stats}) = 0.087
    \]

    \column{0.5\textwidth}
    
  \end{columns}


}
\only<2>{
\begin{columns}
    \column{0.5\textwidth}
    \textbf{Original Data (Counts)}

    \vspace{0.2cm}
    \begin{center}
    \begin{tabular}{ccc}
      \hline
      Math & Stats & CS \\
      \hline
      48  & 230 & 59  \\
      286 & 181 & 73  \\
      286 & 98  & 147 \\
      35  & 14  & 12  \\
      255 & 104 & 235 \\
      $\vdots$ & $\vdots$ & $\vdots$ \\
      \hline
    \end{tabular}
    \end{center}

    \vspace{0.2cm}
    \[
      \mathrm{corr}(\text{Math},\,\text{Stats}) = 0.087
    \]

    \column{0.5\textwidth}
    \textbf{Compositional Data (Normalized to 100)}

    \vspace{0.2cm}
    \begin{center}
    \begin{tabular}{ccc}
      \hline
      Math & Stats & CS \\
      \hline
      14.2 & 68.2 & 17.5 \\
      53.0 & 33.5 & 13.5 \\
      53.9 & 18.5 & 27.7 \\
      57.4 & 23.0 & 19.7 \\
      42.9 & 17.5 & 39.6 \\
      $\vdots$ & $\vdots$ & $\vdots$ \\
      \hline
    \end{tabular}
    \end{center}

    \vspace{0.2cm}
    \[
      \mathrm{corr}(X_{\text{math}},\, X_{\text{stats}}) = -0.874
    \]
  \end{columns}

}

\end{frame}

\end{frame}
\begin{frame}{Spurious Correlation}

 


\only<1>{
\note{Remove this slide:
Spurious correlation

One important consequence of working with compositional data is spurious correlation.

This phenomenon was introduced by Pearson in 1897.

He showed that ratios that share a common denominator can appear correlated, even when the original variables are completely independent.

To see this, let us consider three independent random variables: A, B, and C.

Now suppose we construct two new variables by dividing A by C and B by C.

Even though A and B are independent, these two ratios can become correlated because they share the same denominator.

}
  \begin{block}{Pearson (1897)}
    Ratios that share a common denominator are correlated even when the
    original variables are \emph{completely independent}.
  \end{block}

  \vspace{0.4em}
  \textbf{Simple illustration:} Let $A$, $B$, $C$ be mutually independent. Then:
  \[
    \operatorname{corr}\!\left(\frac{A}{C},\;\frac{B}{C}\right) \neq 0
  \]
  \vspace{0.2em}
  The shared denominator $C$ couples the two ratios artificially.
  The observed correlation reflects the \textbf{denominator's variability},
  not any true relationship between $A$ and $B$.
}

\only<2>{
\note{Spurious correlation

One important consequence of working with compositional data is spurious correlation.

This phenomenon was introduced by Pearson in 1897.

He showed that ratios that share a common denominator can appear correlated, even when the original variables are completely independent.

To see this, let us consider three independent random variables: A, B, and C.

Now suppose we construct two new variables by dividing A by C and B by C.

Even though A and B are independent, these two ratios can become correlated because they share the same denominator.

We observe exactly the same phenomenon in compositional data.

Since all components are normalized by the same sum, a common dependency is introduced.

As a result, even if two components are independent in the original space, they may appear dependent in the simplex space.

However, this is not a mistake.

}
  \begin{block}{Pearson (1897)}
    Ratios that share a common denominator are correlated even when the
    original variables are \emph{completely independent}.
  \end{block}

  \vspace{0.4em}

  \begin{columns}[t]
    \begin{column}{0.38\textwidth}
      \textbf{Original vector:}\\[4pt]
      $\mathbf{Y} = (Y_1, Y_2, \ldots, Y_K) \in \mathbb{R}_+^K$
    \end{column}
    \begin{column}{0.58\textwidth}
      \textbf{Simplex (normalized):}\\[4pt]
      $\mathbf{s}(\mathbf{Y}) = \!\left(\dfrac{Y_1}{S},\, \dfrac{Y_2}{S},\, \ldots,\, \dfrac{Y_K}{S}\right),
      \quad S = \sum_{j=1}^{K} Y_j$
    \end{column}
  \end{columns}

  \vspace{0.5em}
  \rule{\textwidth}{0.3pt}
  \vspace{0.3em}

  Even if $\operatorname{corr}(Y_1,\, Y_2) = 0$, we usually observe:
  \[
    \boxed{\operatorname{corr}\!\left(\frac{Y_1}{S},\,\frac{Y_2}{S}\right) \neq 0}
  \]
}

\only<3>{
\note{Mathematically, everything is correct.

The problem is only interpretation.

Pearson’s warning was exactly about this point: we should not interpret correlations between ratios as correlations between the original variables.

Unfortunately, this mistake often happens in practice.

}
  \[
  \boxed{
    \operatorname{corr}\!\left(\frac{Y_1}{S}, \frac{Y_2}{S}\right)
    =
    \frac{
      \mathbb{E}\!\left(\frac{Y_1 Y_2}{S^2}\right)
      - \mathbb{E}\!\left(\frac{Y_1}{S}\right)\mathbb{E}\!\left(\frac{Y_2}{S}\right)
    }{
      \sqrt{\operatorname{Var}\!\left(\frac{Y_1}{S}\right)}\;
      \sqrt{\operatorname{Var}\!\left(\frac{Y_2}{S}\right)}
    }
  }
  \]
  \vspace{0.4em}
  \begin{block}{Pearson's warning}
    The danger arises when one interprets a correlation between ratios as evidence
    of a genuine association between the \emph{original} variables.
    Such an inference is \textbf{not justified} -- the correlation may exist
    entirely because of the shared denominator, with no signal from the data at all.
  \end{block}
}



\only<4>{
  \note{
  For example, consider a job application dataset with different academic backgrounds such as mathematics, statistics, and computer science.

In the original count data (On the left), the correlation between mathematics and statistics may be very small, close to zero.

However, after normalization to percentages — that is, after moving to the simplex space — the same variables can become strongly negatively or positively correlated.

      }

  \begin{columns}
    \column{0.5\textwidth}
    \textbf{Original Data (Counts)}

    \vspace{0.2cm}
    \begin{center}
    \begin{tabular}{ccc}
      \hline
      Math & Stats & CS \\
      \hline
      48  & 230 & 59  \\
      286 & 181 & 73  \\
      286 & 98  & 147 \\
      35  & 14  & 12  \\
      255 & 104 & 235 \\
      $\vdots$ & $\vdots$ & $\vdots$ \\
      \hline
    \end{tabular}
    \end{center}

    \vspace{0.2cm}
    \[
      \mathrm{corr}(\text{Math},\,\text{Stats}) = 0.087
    \]

    \column{0.5\textwidth}
    \textbf{Compositional Data (Normalized to 100)}

    \vspace{0.2cm}
    \begin{center}
    \begin{tabular}{ccc}
      \hline
      Math & Stats & CS \\
      \hline
      14.2 & 68.2 & 17.5 \\
      53.0 & 33.5 & 13.5 \\
      53.9 & 18.5 & 27.7 \\
      57.4 & 23.0 & 19.7 \\
      42.9 & 17.5 & 39.6 \\
      $\vdots$ & $\vdots$ & $\vdots$ \\
      \hline
    \end{tabular}
    \end{center}

    \vspace{0.2cm}
    \[
      \mathrm{corr}(X_{\text{math}},\, X_{\text{stats}}) = -0.874
    \]
  \end{columns}

  \vspace{0.3em}
  \textbf{Key message:} Spurious dependence:
  This correlation does \textbf{not} necessarily reflect a true relationship.
}

\end{frame}


% ------------------------------------------------------------------
%  SLIDE 2 – Statistical Difficulties: Singular Matrix
% ------------------------------------------------------------------
\begin{frame}{Statistical Difficulties}

\note{
Singularity of the covariance matrix

Another important consequence is that the covariance matrix is not invertible in the simplex space.

This can be shown quite easily.

Because of the constant-sum constraint, the components are not independent.

As a result, one degree of freedom is lost, and the covariance matrix becomes singular.

This creates a practical limitation.

Many standard statistical methods require the inverse of the covariance matrix, such as regression models and multivariate techniques.

Therefore, these methods cannot be directly applied to compositional data.

There are also several other difficulties that arise from working in the simplex space.

I will not go into all of them here.

However, all of these challenges motivated researchers to develop new methods specifically designed for compositional data.
}

\only<1>{
  \vspace{0.4em}
  \[
  \Sigma_{\mathbf{X}} =
  \begin{pmatrix}
    \operatorname{Var}(X_1)     & \operatorname{cov}(X_1,X_2) & \cdots & \operatorname{cov}(X_1,X_K) \\
    \operatorname{cov}(X_2,X_1) & \operatorname{Var}(X_2)     & \cdots & \operatorname{cov}(X_2,X_K) \\
    \vdots                      & \vdots                       & \ddots & \vdots \\
    \operatorname{cov}(X_K,X_1) & \operatorname{cov}(X_K,X_2) & \cdots & \operatorname{Var}(X_K)
  \end{pmatrix}
  \in \mathbb{R}^{K \times K}
  \]
  \vspace{0.2em}
  The \textbf{last row (and column) is linearly determined} by the others via the
  sum-to-constant constraint, confirming $\operatorname{rank}(\Sigma_{\mathbf{X}}) = K-1$.
}

\only<2>{
  The sum-to-constant constraint causes the \textbf{sample covariance matrix to be singular}:
  \[
    \sum_{i=1}^{K} \operatorname{cov}(X_1, X_i)
    = \operatorname{cov}\!\left(X_1,\, \sum_{i=1}^K X_i\right)
    = \operatorname{cov}(X_1,\, 1) = 0
  \]
  \[
    \Rightarrow\quad
    \sum_{i=2}^{K} \operatorname{cov}(X_1, X_i) = -\operatorname{Var}(X_1).
  \]

  \vspace{0.3em}
  $\Sigma_{\mathbf{X}}$ has rank $K-1$ $\;\Rightarrow\;$ \textbf{not invertible}.

  \vspace{0.5em}
  \begin{columns}[T]
    \begin{column}{0.48\textwidth}
      \textbf{Technical consequence:}
      \begin{itemize}
        \item Regression, LDA, and standard PCA all require $\Sigma^{-1}$ -- unavailable.
      \end{itemize}
    \end{column}
    \begin{column}{0.48\textwidth}
      \textbf{Conceptual consequence:}
      \begin{itemize}
        \item No universally accepted definition of correlation for compositional data exists yet.
      \end{itemize}
    \end{column}
  \end{columns}
}

\end{frame}


% ------------------------------------------------------------------
%  SLIDE 3 – The Log-Ratio Idea
% ------------------------------------------------------------------
